\subsection{INTEGRAL OF A LIMIT OF FUNCTIONS ON A COMPACT INTERVAL}

Th. 1 of II, p.~52, applied to the particular case of regulated functions on a compact interval, translates as follows into the notation appropriate to integrals:

PROPOSITION 1. Let A be a set filtered by a filter \(\mathfrak{F}\) , and \((\mathbf{f}_{\alpha})_{\alpha \in \mathrm{A}}\) a family of regulated functions on a compact interval \(\mathrm{I} = [a, b]\) ; if the functions \(\mathbf{f}_{\alpha}\) converge uniformly on I to a (regulated) function \(\mathbf{f}\) with respect to the filter \(\mathfrak{F}\) , then

\[
\lim _ {\mathfrak {F}} \int_ {a} ^ {b} \mathbf {f} _ {\alpha} (t) d t = \int_ {a} ^ {b} \mathbf {f} (t) d t.\tag{1}
\]

Two corollaries to this proposition are important in applications:

COROLLARY 1. Let \((\mathbf{f}_n)\) be a sequence of regulated functions on a compact interval \(\mathrm{I} = [a, b]\) . If the sequence \((\mathbf{f}_n)\) converges uniformly on \(\mathrm{I}\) to a (regulated) function \(\mathbf{f}\) , one has

\[
\lim _ {n \rightarrow \infty} \int_ {a} ^ {b} \mathbf {f} _ {n} (t) d t = \int_ {a} ^ {b} \mathbf {f} (t) d t.\tag{2}
\]

In particular, if a series whose general term \(u_{n}\) is a regulated function on I, converges uniformly to f on I, then the series with general term \(\int_{a}^{b}\mathbf{u}_{n}(t)dt\) is convergent and its sum is \(\int_{a}^{b}\mathbf{f}(t)dt\) (``term-by-term integration of a uniformly convergent series'').

COROLLARY 2. Let A be a subset of a topological space F, and f a map from I × A into a complete normed space E over R, such that, for each \(\alpha \in A\) , the function \(x \mapsto \mathbf{f}(x, \alpha)\) is regulated on I. If the functions \(x \mapsto \mathbf{f}(x, \alpha)\) converge uniformly on I to a (regulated) function \(x \mapsto \mathbf{g}(x)\) , as \(\alpha\) tends to a point \(\alpha_{0} \in \overline{A}\) while remaining in A, then one has

\[
\lim _ {\alpha \rightarrow \alpha_ {0}, \alpha \in \mathrm{A}} \int_ {a} ^ {b} \mathbf {f} (x, \alpha) d x = \int_ {a} ^ {b} \mathbf {g} (x) d x.\tag{3}
\]

In particular:

PROPOSITION 2 (``continuity of an integral with respect to a parameter''). Let F be a compact space, let I = {[}a, b{]} be a compact interval in R, and let f be a continuous map of I × F into a complete normed space E over R; then the function \(\mathbf{h}(\alpha) = \int_{a}^{b} \mathbf{f}(x, \alpha) dx\) is continuous on F.

Indeed, since \(\mathbf{f}\) is uniformly continuous on the compact space \(\mathrm{I} \times \mathrm{F}\) , the functions \(\mathbf{f}(x, \alpha)\) converge uniformly to \(\mathbf{f}(x, \alpha_0)\) on \(\mathrm{I}\) , when \(\alpha\) tends to an arbitrary point \(\alpha_0 \in \mathrm{F}\) .

Here is an application of this proposition: the function \((x, \alpha) \mapsto x^{\alpha}\) is continuous on the product I × J, where I = {[}a, b{]} is a compact interval such that 0 \textless{} a \textless{} b, and J is any compact interval in R; one concludes that \(\int_{a}^{b} x^{\alpha} dx\) is a continuous function of \(\alpha\) on R; now, for \(\alpha\) rational and \(\neq -1\) , this function is equal to \(\frac{b^{\alpha+1} - a^{\alpha+1}}{\alpha + 1}\) , and the function \(\alpha \mapsto \frac{b^{\alpha+1} - a^{\alpha+1}}{\alpha + 1}\) is continuous on every interval of R not containing -1; one thus has (extension of identities) \(\int_{a}^{b} x^{\alpha} dx = \frac{b^{\alpha+1} - a^{\alpha+1}}{\alpha + 1}\) for all real \(\alpha \neq -1\) ; this again means that, for all real \(\alpha\) , the derivative of \(x^{\alpha}\) is \(\alpha x^{\alpha-1}\) (cf.~III, p.~94).

